\documentclass[10pt,a4paper]{amsart}
\usepackage{amsmath,amssymb,amsthm,graphicx}
% Minimal class-compatibility preamble for a reader excerpt of
%   Bhimarthi Ravinder, "Demazure Modules, Chari-Venkatesh Modules and Fusion Products",
%   SIGMA 10 (2014), 110, DOI 10.3842/SIGMA.2014.110.
%
% The author's pinned TeX file declares no theorem environments of its own: they are
% supplied by the SIGMA publisher class sigma.cls, which is not installed here.  The
% block below is that class's theorem-environment block (same environment names and
% styles), copied mechanically so that the retained author text compiles unchanged.
% No mathematical macro or statement is introduced by this file.
\newtheorem{theorem}{Theorem}
\newtheorem*{theorem*}{Theorem}
\newtheorem{lemma}{Lemma}
\newtheorem*{lemma*}{Lemma}
\newtheorem{corollary}{Corollary}
\newtheorem*{corollary*}{Corollary}
\newtheorem{proposition}{Proposition}
\newtheorem*{proposition*}{Proposition}
\newtheorem{conjecture}{Conjecture}
\newtheorem*{conjecture*}{Conjecture}
{\theoremstyle{definition} \newtheorem{definition}{Definition}
\newtheorem*{definition*}{Definition}
\newtheorem{example}{Example}
\newtheorem*{example*}{Example}
\newtheorem{remark}{Remark}
\newtheorem*{remark*}{Remark}
\newtheorem{note}{Note}
\newtheorem*{note*}{Note}
}
\usepackage{cleveref}
\crefname{originalreference}{original reference}{original references}
\Crefname{originalreference}{Original reference}{Original references}
\makeatletter
\newlabel{dm1}{{dm1}{source}{}{sourcepointer-1}{}}
\newlabel{dm1@cref}{{[originalreference][0][]dm1}{[1][0][]source}{}{}{}}
\makeatother
\crefname{originalreference}{original reference}{original references}
\Crefname{originalreference}{Original reference}{Original references}
\makeatletter
\newlabel{dm2}{{dm2}{source}{}{sourcepointer-2}{}}
\newlabel{dm2@cref}{{[originalreference][0][]dm2}{[1][0][]source}{}{}{}}
\makeatother
\title{Explicit defining relations for the level one Demazure module attached to a multiple of the highest root}
\author{Bhimarthi Ravinder}
\date{doi:10.3842/SIGMA.2014.110; CC BY 4.0}
\begin{document}
\maketitle
\noindent\textit{Editorial note.} All delivered mathematical text is the author's own, copied from the pinned article source sources/sigma-2014-110-demazure-fusion/pinned-original-source.tex; no statement or proof step was rewritten, reordered or removed.\par\medskip
\noindent\textit{Editorial note.} Preamble provenance: the author's tex declares no theorem environments of its own (they come from the SIGMA publisher class); the staged sources/sigma-2014-110-demazure-fusion/retained-preamble.tex reproduces that class's theorem-environment block verbatim and is used as the preamble, with the wrapper class and base packages supplied by the assembler.\par\medskip
\noindent\textit{Editorial note.} Bibliography provenance: neither the counted statement nor its proof contains a citation (this lane rescans the whole span before writing the card), so the card carries no bibliography and no entry was reconstructed.\par\medskip
\noindent\textit{Editorial note.} Selection provenance: this card was opened under the file-based control contract (CONTROL.md), with the mechanical pre-card gate run on its statement and proof interval before the card was written.\par\medskip
\noindent\textit{Allowed standard prerequisites (editorial).}\begin{itemize}
\item The level one Demazure module D(1, k theta) as a cyclic graded module over the current algebra with a fixed highest weight vector. The author's section 3, in which the module is set up before the counted statement. Fixes the module whose presentation by generators and relations is asserted.
\item The relations labelled dm1 and dm2 in the article, describing the action of the negative root vectors on the highest weight vector of a Demazure module, and the relation labelled w2'. All three are referred to inside the counted proof (source lines 430, 440) as navigation-pointer stubs. Allow the proof to compare the listed relations with the known ones, in particular to see which of them do not occur.
\item The explicit defining relations of the local Weyl module W(k theta). Recalled at the end of the counted proof (source line 441). Identifies the remaining relations of the statement with the known presentation, which concludes the proof.
\end{itemize}\medskip
\ifdefined\hypertarget\hypertarget{sourcepointer-1}{}\fi
\noindent\textit{Source reference.} Original reference dm1: sources/sigma-2014-110-demazure-fusion/pinned-original-source.tex, source lines 430--430. Navigation pointer only: the counted text refers to the item dm1 of this same article, which is not part of this card. The pointer carries the item's own author label as its number; its content is not reproduced here.\par\medskip
\ifdefined\hypertarget\hypertarget{sourcepointer-2}{}\fi
\noindent\textit{Source reference.} Original reference dm2: sources/sigma-2014-110-demazure-fusion/pinned-original-source.tex, source lines 430--430. Navigation pointer only: the counted text refers to the item dm2 of this same article, which is not part of this card. The pointer carries the item's own author label as its number; its content is not reproduced here.\par\medskip
\begin{proposition}
\label{WvsD}
Given $k\geq 1$, the level~$1$ Demazure module $D(1,k\theta)$ is the graded $\mathfrak{g}[t]$-module generated by an
element $\overline{w}_{k\theta}$, with the following defining relations:
\begin{gather*}
\mathfrak{n}^{+}[t]\, \overline{w}_{k\theta}=0,
\qquad
(h \otimes t^s) \overline{w}_{k\theta} = \langle k\theta, h\rangle \delta_{s,0} \overline{w}_{k\theta},
\qquad
s\geq0,
\quad
h\in\mathfrak{h},
\\
(x^{-}_{\alpha}\otimes 1) \overline{w}_{k\theta}=0,
\qquad
\alpha\in R^+,
\quad
(\theta|\alpha)=0,
\\
(x^{-}_{\alpha} \otimes 1)^{kd_{\alpha}+1} \overline{w}_{k\theta}=0,
\qquad
\big(x^{-}_{\alpha} \otimes t^k\big) \overline{w}_{k\theta}=0,
\qquad
\alpha\in R^+,
\quad
(\theta|\alpha)=1,
\\
(x^{-}_{\theta} \otimes 1)^{2k+1} \overline{w}_{k\theta}=0.
\end{gather*}
\end{proposition}

\begin{gather}
\label{w2'}
\big(x^{-}_\alpha \otimes t^{\langle \lambda,
\alpha^{\vee} \rangle}\big)  w_\lambda=0,
\qquad
\alpha \in R^{+},
\end{gather}

\begin{proof}
Observe that, from the abstract theory of root systems $(\theta|\alpha)= 0$ or~1,
$\forall\, \alpha \in R^{+}\setminus \{\theta\}$.
This implies that $\langle k\theta, \alpha^{\vee}\rangle = 0$ or $kd_{\alpha}$, $\forall\, \alpha \in R^{+}\setminus \{\theta\}$.
Hence the relations~\eqref{dm2} do not occur in $D(1,k\theta)$ and the relations~\eqref{dm1} are
\begin{gather*}
\big(x^{-}_{\alpha} \otimes t^k\big) \overline{w}_{k\theta}=0,
\qquad
\alpha\in R^+,
\quad
\alpha~~\text{short},
\quad
(\theta|\alpha)=1.
\end{gather*}
For a~long root $\alpha\in R^+$ with $(\theta|\alpha)=1$, by~\eqref{w2'} it follows that $(x^{-}_{\alpha} \otimes t^k) \overline{w}_{k\theta}=0$.
Now the other relations are precisely the def\/ining relations of $W(k\theta)$. This proves Proposition~\ref{WvsD}.
\end{proof}

\end{document}
